When using LLMs for empirical studies in SE, researchers face unique challenges and potential limitations that can influence the validity, reliability, and reproducibility of their findings~\cite{sallou2024breaking}.
It is important to openly discuss these limitations and explain how their impact was mitigated.
All this is relative to the current capabilities of LLMs and the current state of the art in terms of tool architectures.
If and how the performance of LLMs in a particular study context will improve with future model generations or new architectural patterns is beyond what researchers can and should discuss as part of a paper's limitation section.
Nevertheless, risk management and threat mitigation are important tasks during the design of an empirical study.
They should not happen as an afterthought.

\guidelinesubsubsection{Recommendations}

\removed{A cornerstone of open science is the ability to reproduce research results.
Although the inherent non-deterministic nature of LLM is a strength in many use cases, its impact on \emph{reproducibility} is problematic.
To enable reproducibility, researchers \should disclose a replication package for their study.
They \should perform multiple repetitions of their experiments (see \benchmarksmetrics) to account for non-deterministic outputs.
In some cases, researchers \may reduce the output variability by setting the temperature to a value close to 0 and setting a fixed seed value.
However, configuring a lower temperature can negatively impact task performance, and not all models allow the configuration of seed values.
Besides the inherent non-determinism, the behavior of an LLM depends on many external factors such as prompt variations and model evolution.
To ensure reproducibility, researchers \should follow all previous guidelines and further discuss the generalization challenges outlined below.}

Researchers \must make every effort to present the limitations of their work clearly \removed{ and forthrightly}
without defensiveness or obfuscation. These limitations may concern diverse topics including generalizability, internal validity (e.g. data leakage), reliability (i.e. non-determinism), and reproducibility (e.g. resource requirements). 

\added{LLM-based studies may involve many kinds of generalization including the following.
\begin{enumerate}
    \item From tested LLMs to others; that is, do the performance characteristics of the LLM(s) studied generalize to LLM(s) not included in the study.
    \item From tested configurations to others; that is, how sensitive are the results to the specific configuration(s) of the LLM under test. 
    \item From research to practice; for instance, just because researchers can get a certain level of code generation performance under ideal conditions does not mean software developers, given the same tools, will get similar results without having the researchers present to show them exactly what to do. 
    \item From a sample of people (e.g. human validators; human participants) to a larger population of people.
    \item From one time period to another; that is, does the same LLM, or other versions of the same LLM, produce similar results at a different point in time?
\end{enumerate}}

If generalizing to multiple LLMs is out of scope, this \must be clearly explained in the \paper; if generalizability is in scope, researchers \must assess generalizability by comparing their results or subsets of the results (e.g., due to computational cost) with other comparable LLMs (see also Section \openllm). Issues with generalizing across configurations, from research to practice (ecological validity), and across people should be discussed in the paper's limitations section. 

The performance of proprietary LLMs (e.g., GPT) can decrease over time for certain tasks using the same \removed{model version (e.g. GPT-4o)}\added{precise model version}~\cite{DBLP:journals/corr/abs-2307-09009, doi:10.1148/radiol.232411}; therefore, reporting the model version and configuration is insufficient to address time-based generalization. The only known way to mitigate this limitation is using an open LLM with a versioning schema and archiving (see \openllm); therefore, researchers \should employ open LLMs to establish a reproducible baseline and, if the use of an open LLM is not possible, researchers \should assess stability and test-retest reliability by repeating their tests over an extended period.

Data leakage, contamination (e.g. of benchmarks), overfitting 
%\removed{occurs when information outside the training data influences model performance.}
\added{from inter-dataset duplication, and growing reliance on big datasets increasingly undermines generalizability. Researchers \mustnot leak their fine-tuned or evaluation datasets into LLM improvement processes.}
%\removed{With the growing reliance on big datasets, the risks of inter-dataset duplication increases} (see, e.g.,~\cite{DBLP:journals/pacmpl/LopesMMSYZSV17, DBLP:conf/oopsla/Allamanis19}). 
%\removed{In the context of LLMs for SE, this can manifest itself, for example, as training data samples that appear in the fine-tuning or evaluation datasets, potentially compromising the validity of the evaluation results~\cite{DBLP:journals/tse/LopezCSSV25}.}
%\removed{Moreover, ChatGPT's functionality to ``improve the model for everyone'' can result in unintentional data leakage.}
%\removed{Hence, t}\added{T}o ensure the validity of the evaluation results, researchers \should carefully curate the fine-tuning and evaluation datasets to prevent 
%\removed{When publishing the results, researchers \should of course still follow open science practices and publish the datasets as part of their \supplementarymaterial.}
If information about the training data of the employed LLM is available, researchers \should evaluate the inter-dataset duplication \removed{and}\added{. Independent of whether such information is available, researchers} \must discuss potential \removed{data leakage or contamination} \added{inter-dataset duplication} \added{and its impact on the reported results} \removed{in the \paper}.
When training an LLM from scratch, researchers \added{face the issue of inter-dataset duplication in the training dataset, biasing the learning process to often-occurring data}. Hence, researchers \added{\should} consider using open datasets \added{curated with deduplication in mind} such as \emph{RedPajama (Together~AI)}~\cite{together2023redpajama}\added{. This enables informed discussion about inter-dataset duplication and its impacts by researchers training and using the LLM.}%\removed{, which are already built with deduplication in mind (with the positive side effect of potentially improving performance~\cite{DBLP:conf/acl/LeeINZECC22}).}

\added{LLM-based research as a resource-intensive process requires transparent reporting of costs and infrastructure to ensure reproducibility and accessibility.}
%Conducting studies with LLMs is a resource-intensive process.
%For self-hosted LLMs, the respective hardware must be provided, and for managed LLMs, the service \emph{costs} must be considered.
The challenge becomes more pronounced as LLMs grow, research architectures become more complex, and experiments become more computationally expensive.
For example, multiple repetitions to assess model or tool performance (see \benchmarksmetrics) multiply the cost and impact \emph{scalability}.
Consequently, resource-intensive research remains predominantly the domain of private companies or well-funded research institutions, hindering researchers with limited resources in reproducing, replicating, or extending study results.
Hence, for transparency reasons, researchers \should report the cost associated with executing an LLM-based study. 
If the study used self-hosted LLMs, researchers \should report the specific hardware used. 
If the study used managed LLMs, the service cost \should be reported\added{, preferably in terms of the number of input tokens, output tokens, and cost per unit.}
To ensure the validity and reproducibility of the research results, researchers \must provide the LLM outputs as evidence (see Section \prompts).
Depending on the architecture (see Section \toolarchitecture), these outputs \should be reported on different architectural levels (e.g., outputs of individual LLMs in multi-agent systems).
Additionally, researchers \should include a subset of the employed validation dataset, selected using an accepted sampling strategy, to allow partial replication of the results.

\emph{Sensitive data} can range from personal to proprietary data, each with its own set of \emph{ethical concerns}.
As mentioned above, a big threat to proprietary LLMs and sensitive data is its use for model improvement.
Hence, using sensitive data can lead to privacy and intellectual property (IP) violations.
Another threat is the implicit bias of LLMs that could lead to discrimination or unfair treatment of individuals or groups.
There is also the fear that using LLMs for qualitative research might \enq{reinforce dominant paradigms and biases} and \enq{identify, replicate and reinforce dominant language and patterns}~\cite{jowsey2025reject}.
To mitigate these concerns, researchers \should minimize the sensitive data used in their studies, \must follow applicable regulations (e.g., GDPR) and individual processing agreements, \should create a data management plan outlining how the data is handled and protected against leakage and discrimination, and \added{\should}\removed{\must} apply for approval from the ethics committee of their organization \removed{(if required)}. 

%\enquote{The field of AI is currently primarily driven by research that seeks to maximize model accuracy — progress is often used synonymously with improved prediction quality. This endless pursuit of higher accuracy over the decade of AI research has significant implications for computational resource requirements and environmental footprint. To develop AI technologies responsibly, we must achieve competitive model accuracy at a fixed or even reduced computational and environmental cost.}~\cite{DBLP:conf/mlsys/WuRGAAMCBHBGGOM22}

Although \emph{metrics} such as \emph{BLEU} or \emph{ROUGE} are commonly used to evaluate the performance of LLMs (see Section \benchmarksmetrics), they may not capture other relevant SE-specific aspects such as functional correctness or runtime performance of automatically generated source code~\cite{DBLP:conf/nips/LiuXW023}.
%Researchers \should clarify and state all relevant requirements, and employ and report the metrics to measure the satisfaction of the requirements (e.g., test-case success rate) and \should follow the best practices described in \href{/guidelines/use-baselines-benchmarks-and-metrics}{Use Suitable Baselines, Benchmarks, and Metrics}.
Given the high resource demand of LLMs, in addition to traditional metrics such as accuracy, precision, and recall, or more contemporary metrics such as \emph{pass@k}, \emph{resource consumption} has to become a key indicator of performance. % to assess research progress responsibly. 
While research has predominantly focused on energy consumption during the early phases of building LLMs (e.g., data center manufacturing, data acquisition, training), inference, that is LLM usage\removed{, often becomes similarly or} \added{ can be} even more resource-intensive~\cite{DBLP:conf/mlsys/WuRGAAMCBHBGGOM22, DBLP:journals/corr/abs-2410-02950, JIANG2024202, mitu2024hidden}.
Hence, researchers \should aim for lower resource consumption on the model side.
This can be achieved by selecting smaller (e.g., GPT-4o-mini instead of GPT-4o) or newer models, or by employing techniques such as model pruning, quantization, or knowledge distillation~\cite{mitu2024hidden}.
Researchers \added{\should} further reduce resource consumption by restricting the number of queries, input tokens, or output tokens~\cite{mitu2024hidden}, by using different prompt engineering techniques (e.g., zero-shot prompts seem to emit less CO2 than chain-of-thought prompts), or by carefully sampling smaller datasets for fine-tuning and evaluation instead of using large datasets.
Reducing resource consumption involves a trade-off with our recommendation to perform multiple experimental runs to account for non-determinism, as suggested in \benchmarksmetrics.
To report the environmental impact of a study, \removed{researchers \should use} software such as \href{https://github.com/mlco2/codecarbon}{CodeCarbon} or \href{experiment-impact-tracker}{Experiment Impact Tracker} \removed{to track and quantify} \added{can support in tracking and quantifying} the carbon footprint of the study or report an estimate of the carbon footprint through tools such as \href{https://mlco2.github.io/impact/#about}{MLCO2 Impact}.
\removed{They \should detail the LLM version and configuration as described in \modelversion, state the hardware or hosting provider of the model as described in \toolarchitecture.}
Researchers \must justify why LLMs were chosen over existing approaches and discuss the achieved performance in relation to their potentially higher resource consumption.

\guidelinesubsubsection{Example(s)}

An example highlighting the need for caution around \emph{replicability} is the study of \citeauthor{DBLP:conf/sigir-ap/StaudingerKPLH24}~\cite{DBLP:conf/sigir-ap/StaudingerKPLH24} who attempted to replicate an LLM study.
The authors were unable to reproduce the exact results, even though they saw similar trends as in the original study.
%They consider their results as not reliable enough for a systematic review.

To analyze whether the results of proprietary LLMs transfer to open LLMs, \citeauthor{DBLP:conf/sigir-ap/StaudingerKPLH24} benchmarked previous results using GPT-3.5 and GPT4 against Mistral and Zephyr~\cite{DBLP:conf/sigir-ap/StaudingerKPLH24}.
They found that open-source models could not deliver the same performance as proprietary models. %, restricting the effect to certain proprietary models.
This paper also \added{provides an example of \emph{cost} reporting}: \enquote{120 USD in API calls for GPT 3.5 and GPT 4, and 30 USD in API calls for Mistral AI. Thus, the total LLM cost of our reproducibility study was 150 USD}.
\citeauthor{tinnessoftware} is an example of balancing dataset size between the need for manual semantic analysis and computational resource consumption~\cite{tinnessoftware}.

Individual studies have already begun to highlight \removed{the} uncertainty about the \emph{generalizability} of their results in the future. \citeauthor{DBLP:conf/msr/JesseADM23} acknowledge the issue that LLMs evolve over time and that this evolution could impact the study results~\cite{DBLP:conf/msr/JesseADM23}.
\added{Inter-dataset duplication as another threat to generalizability received extensive attention over the years in SE tasks evolving around code (see, e.g., \cite{DBLP:journals/pacmpl/LopesMMSYZSV17, DBLP:conf/oopsla/Allamanis19, DBLP:journals/ese/KarmakarAR23, DBLP:journals/tse/LopezCSSV25}).}
\removed{The issue of inter-dataset duplication has also attracted interest in other disciplines. For example, in biology, Lakiotaki et al. acknowledge and address the overlap between multiple common disease datasets \cite{DBLP:journals/biodb/LakiotakiVTGT18}.}

\added{An example to mitigate inter-dataset duplication, and hence increase generalizability can be found in \citeauthor{DBLP:conf/ease/CoignionQR24} work. In an effort to evaluated the performance of LLMs to produce LeetCode \citeauthor{DBLP:conf/ease/CoignionQR24} mitigated the issue, by only using} LeetCode problems published after 2023-01-01, reducing the likelihood of LLMs having seen those problems before~\cite{DBLP:conf/ease/CoignionQR24}.
Further, they discuss the performance differences of LLMs on different datasets in light of potential inter-dataset duplication.
\citeauthor{zhou2025lessleakbenchinvestigationdataleakage} performed an empirical evaluation of data leakage in 83 software engineering benchmarks~\cite{zhou2025lessleakbenchinvestigationdataleakage}.
Although most benchmarks suffer from minimal leakage, \removed{very} \added{a} few showed a leakage of up to 100\%.
The authors found a high impact of data leakage on the performance evaluation\added{, impacting generalizability.}
A starting point for studies that aim to assess and mitigate inter-dataset duplication are the \href{https://huggingface.co/datasets/tiiuae/falcon-refinedweb}{Falcon LLMs}, for which parts of its training data are available on Hugging Face~\cite{technology_innovation_institute_2023}.
Through this dataset, it is possible to reduce the overlap between the training and evaluation data, improving the \removed{validity}\added{generalizability} of the evaluation results.
A starting point to prevent active data leakage into a LLM improvement process is to ensure that the data is not used to train the model (e.g., via OpenAI's data control functionality)~\cite{DBLP:conf/eacl/BalloccuSLD24}.

Finally, bias can occur in LLM training datasets, resulting in various types of discrimination.
\citeauthor{DBLP:journals/corr/abs-2309-00770} propose metrics to quantify biases in various tasks (e.g., text generation, classification, question answering)~\cite{DBLP:journals/corr/abs-2309-00770}.

\guidelinesubsubsection{Benefits}

%Ensuring \emph{reproducibility} allows for the independent replication and verification of study results.
The \emph{reproduction} or \emph{replication} of study results under similar conditions by different parties greatly increases the validity of the results. % and promotes research transparency.
Independent verification is of particular importance for studies involving LLMs, due to the non-determinism of their outputs and the potential for biases in their training, fine-tuning, and evaluation datasets.
Mitigating threats to \emph{generalizability} of a study through the integration of an open LLM as a baseline, or the reporting of results over an extended period of time can increase the validity, reliability, and replicability of a study's results.
Assessing and mitigating the effects of \emph{inter-dataset duplication} strengthens a study's validity and reliability, as it prevents overly optimistic performance estimates that do not apply to previously unknown samples.
Reporting the \emph{costs} of performing a study not only increases transparency but also provides context for secondary literature. % in setting primary research into perspective.
Providing \emph{replication packages} with LLM output and data samples for partial replicability are paramount steps toward open and inclusive research in light of resource inequality between researchers.
Considering and justifying the usage of LLMs over other approaches can lead to more efficient and sustainable solutions for SE problems. 
Reporting the \emph{environmental impact} of LLM usage also sets the stage for more sustainable research practices in SE.

\guidelinesubsubsection{Challenges}

With commercial LLMs evolving over time, the \emph{generalizability} of results to future model versions is uncertain.
Employing open LLMs as a baseline can mitigate this limitation, but may not always be feasible due to computational cost.
Most LLM providers do not publicly offer information about their training data, impeding the assessment of \emph{inter-dataset duplication} effects.
Consistently keeping track of and reporting the \emph{costs} involved in a research endeavor is challenging.
Building a coherent \emph{replication package} that includes LLM outputs and samples for partial replicability requires additional effort and resources.
Finding the right \emph{metrics} to evaluate the performance of LLMs in SE for specific tasks is challenging.
Defining all requirements beforehand to ensure the use of suitable metrics can be a challenge, especially in exploratory research.
Our Section \benchmarksmetrics and its references can serve as a starting point.
Ensuring \emph{compliance} across jurisdictions is difficult with different regions having different regulations and requirements (e.g., GDPR and the AI Act in the EU, CCPA in California).
Selecting datasets and models with fewer \emph{biases} is challenging, as biases are often unknown.
Measuring or estimating the \emph{environmental impact} of a study is challenging and might not always be feasible.
\removed{Especially in exploratory research,} \added{In exploratory research in particular,} the impact is hard to estimate in advance, making it difficult to justify the usage of LLMs over other approaches.

\guidelinesubsubsection{Study Types}

Limitations and mitigations \should be discussed for all study types in relation to the context of the individual study.
This section provides a starting point and lists aspects along which researchers can reflect on the limitations of their study design and potential mitigations that exist.

\guidelinesubsubsection{Advice for Reviewers}

\added{We suggest that reviewers begin by checking the following:
\begin{enumerate}
    \item The paper has \removed{a section or subsection explicitly} \added{an explicit section or sub-section} with a name like ``limitations'' or ``threats to validity''.
    \item The \removed{limitations} section briefly describes numerous limitations concerning multiple dimensions of research rigor (e.g. internal, external, and construct validity, reliability) 
    \item The tone of the \removed{limitations} section is forthright and nondefensive; it does not try to downplay the limitations. 
\end{enumerate}}

\added{A reviewer noticing undisclosed limitations is normal and typically not grounds for rejection. The reviewer should simply ask for these additional limitations to be disclosed. Similarly, reviewers should not penalize authors for long lists of limitations, or studies that seem quite limited. Most research seems very limited when the authors are forthright about its limitations. In contrast, reviewers should be wary of any paper that appears to have few limitations, or emphasizes how minor they are.} 

\added{Reviewers should encourage authors to structure their limitation section according to well-understood categories rather than over-using LLM-specific jargon; for example, data leakage and contamination are \textit{internal validity} concerns while non-determinism causes \textit{test-retest reliability} concerns. LLM-specific limitations should be connected to well-understood scientific quality criteria, not presented as fundamentally new and different.} 

\added{Reviewers may find \remove{assessing limitations hindered by} \added{assessing limitations difficult due to} disorganized and confused limitations sections. % Brian: I think this is unnecessary. No point in taking a general swipe at the SE research community here:        Many SE researchers do not understand the vocabulary of research quality. This is why papers frequently discuss the ``internal validity'' of studies with no causal relationships, the ``construct validity'' of studies with no constructs, or the ``objectivity'' of qualitative research that is intentionally and necessarily subjective.  Reviewers have a moral and professional duty to familiarize themselves with the language of science. If you cannot clearly explain the different types of validity, avail yourself of a good book on research methods (e.g.~\cite{trochim2001research}) and encourage others to do the same. 